\chapter{Functorial Semantics: Rendering, Transport, and Synthesis}

A unified project should not reserve its mathematics for only one corner of the system. If
local-to-global semantics is truly the repository's common language, then even apparently
secondary concerns such as rendering, transport, and synthesis should become legible in that
language. This chapter argues that they do. The right conceptual move is to treat them as
functorial constructions on typed presheaves rather than as utility layers sitting outside
the theory.

\section*{Rendering as a natural transformation}
Suppose \(\mathcal{F}\) is a presheaf of typed local evidence. A renderer does not discover a
new semantic fact ex nihilo. It reorganizes the existing typed content into another target
language: English explanation, Lean code, diagrams, tables, or machine-readable reports.
That suggests the right abstraction. Rendering is a natural transformation from a presheaf of
typed sections into a presheaf of representations.

In the repository, this idea appears in multiple places. The web-facing explanations in
\repo{web/} are not external to the type story; they are human-readable images of it. The
Lean translation pipeline in \repo{src/deppy/hybrid/lean_translation/} is not just export; it
is a semantics-preserving map from one typed universe into another. The report and rendering
machinery in \repo{src/deppy/render/} likewise maps trust-bearing objects into surfaces that a
user or auditor can inspect.

What matters is that a good renderer respects restriction maps. If one moves from a whole
function to a branch site or from a module to an overlap interface, the rendered object
should continue to agree with the underlying semantics. That is exactly the sort of square
one expects of a natural transformation.

\[
\begin{tikzcd}[column sep=large, row sep=large]
\mathcal{F}(V) \arrow[r, "R_V"] \arrow[d, "\mathrm{res}_{U,V}"'] & \mathcal{R}(V) \arrow[d, "\mathrm{res}_{U,V}"] \\
\mathcal{F}(U) \arrow[r, "R_U"] & \mathcal{R}(U)
\end{tikzcd}
\]

Here \(\mathcal{R}\) might be a Lean-syntax presheaf, a textual-report presheaf, or a diagram
presheaf. The point is the same in every case: rendering should commute with locality.

\section*{Transport as restriction and extension}
Interprocedural analysis and contract propagation are often discussed as if they were
orthogonal to typing. In this project they are not. They are the mechanics by which typed
information moves along a program's semantic edges. One may call that restriction,
extension, transport, or pullback depending on which aspect one wants to emphasize, but it
is still the same local-to-global geometry.

When a caller invokes a callee, one restricts attention from the caller's current state to
the callee's input interface. When the callee returns, one transports postconditions back
into the caller's context. When a wrapper preserves an invariant, one is witnessing a
particular compatibility between the wrapped and unwrapped sites. When a zero-change analyzer
infers a likely precondition from existing source, it is harvesting transportable local
facts from ordinary Python and lifting them into typed form.

This is why the repository's adoption story remains coherent. The same semantics that govern
explicitly annotated code can also govern inferred local structure. The difference is not one
of basic ontology but of how much evidence is given up front.

\section*{Synthesis as cover elaboration}
Perhaps the strongest test of unity is code synthesis. If synthesis sits outside the theory,
then the project fractures: mathematics on one side, generation on the other. The repository
is most interesting precisely because it refuses that fracture. The generation agent is best
understood as an elaborator of covers.

A user prompt does not immediately produce a monolithic program. It produces intent-level
material from which a cover is elaborated: modules, interfaces, invariants, proof targets,
and gluing obligations. Local generation then attempts to produce sections on each chart.
Checking passes test whether those sections agree on overlaps. Residual failures trigger
cover refinement or local repair. Only at the end does one speak about a global artifact.

This account is far superior to saying that the repository simply ``uses an LLM and then
verifies the output.'' That description misses the architecture. The architecture says that
synthesis itself is a typed local-to-global process.

\begin{figure}[h]
\centering
\begin{tikzpicture}[
  node distance=1.0cm,
  box/.style={draw=chapterblue, rounded corners=2pt, fill=chapterblue!6, minimum width=2.8cm, minimum height=0.85cm, align=center},
  obs/.style={draw=accentgold, rounded corners=2pt, fill=accentgold!12, minimum width=3.0cm, minimum height=0.85cm, align=center},
  >=Latex
]
\node[box] (prompt) {prompt};
\node[box, right=of prompt] (cover) {cover elaboration};
\node[box, right=of cover] (sections) {local sections};
\node[obs, right=of sections] (glue) {global artifact or obstruction};
\node[obs, below=of sections] (repair) {cover refinement / local repair};
\draw[->, thick, chapterblue] (prompt) -- (cover);
\draw[->, thick, chapterblue] (cover) -- (sections);
\draw[->, thick, chapterblue] (sections) -- (glue);
\draw[->, thick, accentgold] (glue) -- node[right] {if non-gluing} (repair);
\draw[->, thick, accentgold] (repair.west) .. controls +(-3.0,0.0) and +(0.0,-0.8) .. (cover.south);
\end{tikzpicture}
\caption{Synthesis as cover elaboration. The agent is most coherently understood as a machine for growing, testing, and refining covers until local sections glue.}
\end{figure}

\section*{Lean translation as a deep embedding}
The Lean pipeline deserves special mention because it dramatizes the project's ambition. A
lightweight renderer merely paraphrases typed content. Lean translation aims higher. It tries
to embed a portion of the repository's typed world into a theorem-proving environment whose
kernel can then check part of the resulting content. This is not a casual format conversion.
It is a deep embedding of selected structural and proof-bearing aspects of the stack.

That also explains why residual honesty matters so much. A functor into Lean cannot simply
pretend that every semantic predicate has become a theorem. Some content becomes theorem
statements. Some becomes lemmas with external evidence. Some becomes axioms. Some remains as
residual assumptions. A faithful deep embedding preserves these distinctions rather than
obscuring them.

\section*{Why this strengthens the sense of one project}
The unifying power of the present viewpoint is easiest to appreciate negatively. Without it,
rendering looks like documentation, transport looks like static-analysis plumbing, and
synthesis looks like an attached generation system. With it, they become further instances of
a single mathematical habit: typed local material should move coherently between sites,
between layers, and between external representations.

That is exactly the kind of broad conceptual reuse a monograph needs if it is to feel like a
single project rather than a collection of adjacent papers. The project is not merely about
which types exist. It is about how typed local evidence circulates, is rendered, is upgraded,
is repaired, and is finally assembled.

\section*{The repository as a functorial workshop}
A final image may help. The repository is not only a checker and not only a generator. It is
a workshop of functors and natural transformations built around typed presheaves. Some maps
go downward into code. Some go upward into proofs. Some go sideways into reports, diagrams,
and English explanations. Some shuttle information across call graphs and module boundaries.
What makes these maps members of one family is that they are all expected to respect the same
local-to-global semantics.

Once that expectation is made explicit, the repository reads more like a foundational
research program. It is building not just tools, but a universe in which software artifacts,
proof artifacts, semantic judgments, and generated designs can all be transported through a
single coherent language.
